% formal/ppcmem.tex
% mainfile: ../perfbook.tex
% SPDX-License-Identifier: CC-BY-SA-3.0

\section{专用状态空间搜索}
\label{sec:formal:Special-Purpose State-Space Search}
%
\epigraph{样样通，样样松。}{佚名}

虽然 Promela 和 Spin 允许你验证几乎任何（较小的）算法，
但它们的通用性有时会成为一种负担。
例如，Promela 不理解 \IX{memory model} 或任何类型的
重排序语义。
因此，本节描述了一些理解生产系统使用的 memory model 的
state-space search 工具，大大简化了弱序代码的验证。

例如，
\cref{sec:formal:Promela Example: QRCU}
展示了如何说服 Promela 来考虑弱内存排序。
虽然这种方法可以很好地工作，但它要求开发者
完全理解系统的 memory model。
不幸的是，很少有（如果有的话）开发者完全理解
现代 CPU 的复杂 memory model。

因此，另一种方法是使用已经理解内存排序的工具，
例如由剑桥大学的 \ppl{Peter}{Sewell} 和 \ppl{Susmit}{Sarkar}、
INRIA 的 \ppl{Luc}{Maranget}、\ppl{Francesco Zappa}{Nardelli}
和 \ppl{Pankaj}{Pawan}、牛津大学的 \ppl{Jade}{Alglave}
与 IBM 的 \ppl{Derek}{Williams} 合作产生的
PPCMEM 工具~\cite{JadeAlglave2011ppcmem}。
这个团队形式化了 Power、\ARM、x86 以及
C/C++11 标准~\cite{RichardSmith2019N4800}的 memory model，
并基于 Power 和 \ARM\ 形式化产生了 PPCMEM 工具。

\QuickQuiz{
	但 x86 有强内存排序，为什么要形式化它的
	memory model？
}\QuickQuizAnswer{
	实际上，学术界认为 x86 的 memory model 是弱的，
	因为它可以允许先前的存储与后续的加载重排序。
	从学术角度来看，强 memory model 是完全不允许
	重排序的，这样所有线程都同意它们可见的所有操作的顺序。

	而且确实存在开发者有时对 x86 内存排序
	感到困惑的情况。
}\QuickQuizEnd

PPCMEM 工具以 \emph{litmus test} 作为输入。
一个示例 litmus test 在
\cref{sec:formal:Anatomy of a Litmus Test}中介绍。
\Cref{sec:formal:What Does This Litmus Test Mean?}
将这个 litmus test 与等效的 C 语言程序联系起来，
\cref{sec:formal:Running a Litmus Test}描述了如何
将 PPCMEM 应用于这个 litmus test，而
\cref{sec:formal:PPCMEM Discussion}
讨论了其影响。

\subsection{Litmus Test 的剖析}
\label{sec:formal:Anatomy of a Litmus Test}

\cref{lst:formal:PPCMEM Litmus Test}展示了
PPCMEM 的一个 PowerPC litmus test 示例。
ARM 接口以相同方式工作，但用 \ARM\ 指令
替代 Power 指令，并将开头的 \qco{PPC}
替换为 \qco{ARM}。

\begin{listing}
\begin{fcvlabel}[ln:formal:PPCMEM Litmus Test]
\begin{VerbatimL}[commandchars=\@\[\]]
PPC SB+lwsync-RMW-lwsync+isync-simple		@lnlbl[type]
""						@lnlbl[altname]
{						@lnlbl[init:b]
0:r2=x; 0:r3=2; 0:r4=y; 0:r10=0; 0:r11=0; 0:r12=z; @lnlbl[init:0]
1:r2=y; 1:r4=x;					@lnlbl[init:1]
}						@lnlbl[init:e]
 P0                 | P1           ;		@lnlbl[procid]
 li r1,1            | li r1,1      ;		@lnlbl[reginit]
 stw r1,0(r2)       | stw r1,0(r2) ;		@lnlbl[stw]
 lwsync             | sync         ; @lnlbl[P0lwsync] @lnlbl[P1sync]
                    | lwz r3,0(r4) ; @lnlbl[P0empty]  @lnlbl[P1lwz]
 lwarx  r11,r10,r12 | ;		@lnlbl[P0lwarx] @lnlbl[P1empty:b]
 stwcx. r11,r10,r12 | ;		@lnlbl[P0stwcx]
 bne Fail1          | ;		@lnlbl[P0bne]
 isync              | ;		@lnlbl[P0isync]
 lwz r3,0(r4)       | ;		@lnlbl[P0lwz]
 Fail1:             | ;		@lnlbl[P0fail1] @lnlbl[P1empty:e]

exists						@lnlbl[assert:b]
(0:r3=0 /\ 1:r3=0)				@lnlbl[assert:e]
\end{VerbatimL}
\end{fcvlabel}
\caption{PPCMEM Litmus Test}
\label{lst:formal:PPCMEM Litmus Test}
\end{listing}

\begin{fcvref}[ln:formal:PPCMEM Litmus Test]
在这个例子中，\clnref{type}标识了系统类型（\qco{ARM} 或
\qco{PPC}）并包含模型的标题。
\Clnref{altname}提供了一个测试的替代名称位置，
通常你会像上面的例子那样留空。
可以使用 OCaml（或 Pascal）的 \nbco{(* *)} 语法
在\clnref{altname,init:b}之间插入注释。

\Clnrefrange{init:b}{init:e}给出所有寄存器的初始值；
每个的形式为
\co{P:R=V}，其中 \co{P} 是进程标识符，\co{R} 是寄存器
标识符，\co{V} 是值。
例如，进程~0 的寄存器 \co{r3} 初始包含值~2。
如果值是一个变量（在例子中是 \co{x}、\co{y} 或 \co{z}），
那么寄存器被初始化为该变量的地址。
也可以初始化变量的内容，例如，
\co{x=1} 将 \co{x} 的值初始化为~1。
未初始化的变量默认为零值，所以在
例子中，\co{x}、\co{y} 和 \co{z} 初始都为零。

\Clnref{procid}提供了两个进程的标识符，所以
\clnref{init:0}上的 \co{0:r3=2} 也可以写成
\co{P0:r3=2}。
\Clnref{procid}是必须的，标识符必须是
\co{Pn} 的形式，其中 \co{n} 是列号，从最左列的零开始。
这可能看起来不必要地严格，但它确实防止了
实际使用中相当多的混淆。
\end{fcvref}

\QuickQuiz{
	\begin{fcvref}[ln:formal:PPCMEM Litmus Test]
	为什么\cref{lst:formal:PPCMEM Litmus Test}的
	\clnref{reginit}要初始化寄存器？
	为什么不在\clnref{init:0,init:1}上初始化它们？
	\end{fcvref}
}\QuickQuizAnswer{
	两种方式都可以。
	然而，一般来说，使用初始化比使用显式指令更好。
	在这个例子中使用显式指令是为了演示它们的使用。
	此外，该工具网站上提供的许多 litmus test
	（\url{https://www.cl.cam.ac.uk/~pes20/ppcmem/}）
	是自动生成的，这会生成显式的初始化指令。
}\QuickQuizEnd

\begin{fcvref}[ln:formal:PPCMEM Litmus Test]
\Clnrefrange{reginit}{P0fail1}是每个进程的代码行。
一个给定的进程可以有空行，如 P0 的
\clnref{P0empty}和 P1 的\clnrefrange{P1empty:b}{P1empty:e}。
标签和分支是允许的，如\clnref{P0bne}上的分支
到\clnref{P0fail1}上的标签所演示的那样。
话虽如此，过于自由地使用分支会扩展 state space。
使用循环是一种特别好的方式来让你的 state space 爆炸。

\Clnrefrange{assert:b}{assert:e}展示了 assertion，在这种情况下
表示我们关心的是在两个线程完成执行后，P0 和 P1 的 \co{r3} 寄存器
是否都能包含零。
这个 assertion 很重要，因为有许多用例
如果 P0 和 P1 在各自的 \co{r3} 寄存器中都看到零，
将会惨败。

这应该给你足够的信息来构造简单的 litmus test。
还有一些额外的文档可用，不过大部分
额外文档是为另一个在实际硬件上运行测试的
研究工具准备的。
也许更重要的是，在线工具提供了大量预先存在的 litmus test
（可通过 \url{https://www.cl.cam.ac.uk/~pes20/ppcmem/}
的 ``Select ARM Test'' 和 ``Select POWER Test'' 按钮获得）。
这些预先存在的 litmus test 中很可能有一个能
回答你的 Power 或 \ARM\ 内存排序问题。

\subsection{这个 Litmus Test 是什么意思？}
\label{sec:formal:What Does This Litmus Test Mean?}

P0 的\clnref{reginit,stw}等价于 C 语句 \co{x=1}，
因为\clnref{init:0}将 P0 的寄存器 \co{r2} 定义为
\co{x} 的地址。
P0 的\clnref{P0lwarx,P0stwcx}分别是 load-linked
（在 \ARM\ 中称为 ``load register exclusive''，
在 Power 中称为 ``load reserve''）
和 store-conditional（在 \ARM\ 中称为
``store register exclusive''）的助记符。
当它们一起使用时，形成一个原子指令序列，
大致类似于以 x86 \co{lock;cmpxchg} 指令为代表的
\IXacrml{cas} 序列。
提升到更高的抽象层次，从
\clnrefrange{P0lwsync}{P0isync}的序列
等价于 Linux 内核的 \co{atomic_add_return(&z, 0)}。
最后，\clnref{P0lwz}大致等价于 C 语句 \co{r3=y}。

P1 的\clnref{reginit,stw}等价于 C 语句 \co{y=1}，
\clnref{P1sync}是一个内存屏障，等价于 Linux 内核语句 \co{smp_mb()}，
\clnref{P1lwz}等价于 C 语句 \co{r3=x}。
\end{fcvref}

\QuickQuiz{
	\begin{fcvref}[ln:formal:PPCMEM Litmus Test]
	但是\cref{lst:formal:PPCMEM Litmus Test}的
	\clnref{P0fail1}，即 \co{Fail1:} 标签那行，
	怎么了？
	\end{fcvref}
}\QuickQuizAnswer{
	PowerPC 版本的 \co{atomic_add_return()} 实现
	在 \co{stwcx} 指令失败时循环，它通过在
	条件码寄存器中设置非零状态来通知这一点，
	然后由 \co{bne} 指令测试。
	因为实际建模循环会导致 state-space 爆炸，
	我们改为跳转到 \co{Fail1:} 标签，
	以 P0 的 \co{r3} 寄存器中的初始值~2 终止模型，
	这不会触发 exists assertion。

	关于这个技巧是否普遍适用存在一些争论，
	但我还没有看到它失败的例子。
}\QuickQuizEnd

\begin{listing}
\begin{VerbatimL}
void P0(void)
{
	int r3;

	x = 1; /* Lines 8 and 9 */
	atomic_add_return(&z, 0); /* Lines 10-15 */
	r3 = y; /* Line 16 */
}

void P1(void)
{
	int r3;

	y = 1; /* Lines 8-9 */
	smp_mb(); /* Line 10 */
	r3 = x; /* Line 11 */
}
\end{VerbatimL}
\caption{PPCMEM Litmus Test 的含义}
\label{lst:formal:Meaning of PPCMEM Litmus Test}
\end{listing}

将所有这些放在一起，整个 litmus test 的 C 语言等价物
如\cref{lst:formal:Meaning of PPCMEM Litmus Test}所示。
关键点是如果 \co{atomic_add_return()} 充当完整的
内存屏障（如 Linux 内核要求的那样），
那么 \co{P0()} 和 \co{P1()} 的 \co{r3}
变量在执行完成后应该不可能都为零。

下一节描述如何运行这个 litmus test。

\subsection{运行 Litmus Test}
\label{sec:formal:Running a Litmus Test}

如前所述，litmus test 可以通过
\url{https://www.cl.cam.ac.uk/~pes20/ppcmem/}交互式运行，
这有助于建立对 \IX{memory model} 的理解。
然而，这种方法要求用户手动进行
完整的 state-space search。
因为很难确保你已经检查了每种可能的事件序列，
所以为此目的提供了一个单独的工具~\cite{PaulEMcKenney2011ppcmem}。

\begin{listing}
\begin{VerbatimL}[numbers=none,xleftmargin=0pt]
./ppcmem -model lwsync_read_block \
         -model coherence_points filename.litmus
...
States 6
0:r3=0; 1:r3=0;
0:r3=0; 1:r3=1;
0:r3=1; 1:r3=0;
0:r3=1; 1:r3=1;
0:r3=2; 1:r3=0;
0:r3=2; 1:r3=1;
Ok
Condition exists (0:r3=0 /\ 1:r3=0)
Hash=e2240ce2072a2610c034ccd4fc964e77
Observation SB+lwsync-RMW-lwsync+isync Sometimes 1
\end{VerbatimL}
\caption{PPCMEM 检测到错误}
\label{lst:formal:PPCMEM Detects an Error}
\end{listing}

因为\cref{lst:formal:PPCMEM Litmus Test}中
展示的 litmus test 包含读-修改-写指令，我们必须在命令行
中添加 \co{-model} 参数。
如果 litmus test 存储在 \co{filename.litmus} 中，
结果将产生如\cref{lst:formal:PPCMEM Detects an Error}所示的输出，
其中 \co{...} 代表大量的进度输出。
状态列表包括 \co{0:r3=0; 1:r3=0;}，再次表明
旧的 PowerPC 实现的 \co{atomic_add_return()} 不能
充当完整屏障。
最后一行的 ``Sometimes'' 确认了这一点：
assertion 在某些执行中触发，但不是每次都触发。

\begin{listing}
\begin{VerbatimL}[numbers=none,xleftmargin=0pt]
./ppcmem -model lwsync_read_block \
         -model coherence_points filename.litmus
...
States 5
0:r3=0; 1:r3=1;
0:r3=1; 1:r3=0;
0:r3=1; 1:r3=1;
0:r3=2; 1:r3=0;
0:r3=2; 1:r3=1;
No (allowed not found)
Condition exists (0:r3=0 /\ 1:r3=0)
Hash=77dd723cda9981248ea4459fcdf6097d
Observation SB+lwsync-RMW-lwsync+sync Never 0 5
\end{VerbatimL}
\caption{修复后的 Litmus Test 上的 PPCMEM}
\label{lst:formal:PPCMEM on Repaired Litmus Test}
\end{listing}

修复这个 Linux 内核 bug 的方法是将 P0 的 \co{isync} 替换为
\co{sync}，这产生了如
\cref{lst:formal:PPCMEM on Repaired Litmus Test}所示的输出。
如你所见，\co{0:r3=0; 1:r3=0;} 没有出现在状态列表中，
最后一行标注了 ``Never''。
因此，模型预测有问题的执行序列不会发生。

\QuickQuizSeries{%
\QuickQuizB{
	\ARM\ Linux 内核有类似的 bug 吗？
}\QuickQuizAnswer{
	\ARM\ 没有这个特定的 bug，因为它在
	\co{atomic_add_return()} 函数的汇编语言实现
	前后都放置了 \co{smp_mb()}。
	PowerPC 也不再有这个 bug；它早已被
	修复~\cite{BenjaminHerrenschmidt2011:powerpc:atomic_return}。
}\QuickQuizEndB
%
\QuickQuizE{
	\begin{fcvref}[ln:formal:PPCMEM Litmus Test]
	\cref{lst:formal:PPCMEM Litmus Test}中
	\clnref{P0lwsync}上的 \co{lwsync} 是否提供了
	足够的排序？
	\end{fcvref}
}\QuickQuizAnswerE{
	这取决于所需的语义。
	本回答的其余部分假设
	\cref{lst:formal:PPCMEM Litmus Test}中
	\co{P0} 的汇编语言旨在实现一个返回值的原子操作。

	正如\cref{chp:Advanced Synchronization: Memory Ordering}中
	所讨论的，
	Linux 内核的内存一致性模型要求
	返回值的原子 RMW 操作在两侧都是完全排序的。
	\co{lwsync} 提供的排序不足以达到此目的，
	因此应该使用 \co{sync} 代替。
	这个更改后来已经
	完成~\cite{BoqunFeng2015:powerpc:value-returning-atomics}，
	以响应讨论其他几个 litmus test 的
	邮件线程~\cite{Paulmck2015:powerpc:value-returning-atomics}。
	找到 Linux 内核可能存在的其他 bug 留作
	读者的练习。

	在提供较弱语义的其他环境中，\co{lwsync}
	可能就够了。
	但对于 Linux 内核的返回值原子操作则不行！
}\QuickQuizEndE
}

\subsection{PPCMEM 讨论}
\label{sec:formal:PPCMEM Discussion}

这些工具有望对从事在 \ARM\ 和 Power 上运行的底层
并行原语工作的人员提供巨大帮助。
这些工具确实有一些固有的局限性：

\begin{enumerate}
\item	这些工具是研究原型，因此不提供支持。
\item	这些工具不构成 IBM 或 \ARM\ 对其各自 CPU 架构的
	官方声明。
	例如，两家公司保留在任何时间对这些工具的任何版本
	报告 bug 的权利。
	因此，这些工具不能替代在真实硬件上进行仔细的
	压力测试。
	此外，工具和它们所基于的模型都在积极开发中，
	可能随时发生变化。
	另一方面，这个模型是与相关硬件专家协商开发的，
	所以有很好的理由相信它是对架构的可靠表示。
\item	这些工具目前处理指令集的一个子集。
	这个子集对我的目的来说已经足够了，但你的情况
	可能不同。
	特别是，该工具只处理字大小的访问（32位），
	并且被访问的字必须正确对齐。\footnote{
		但最近的工作关注混合大小的
		访问~\cite{Flur:2017:MCA:3093333.3009839}。}
	此外，该工具不处理 \ARM\ 内存屏障指令的某些
	较弱变体，也不处理算术运算。
\item	这些工具仅限于在少量线程上运行的小型无循环代码片段。
	更大的例子会导致 state-space 爆炸，
	就像 Promela 和 Spin 等类似工具一样。
\item	完整的 state-space search 不会指出
	每个有问题的状态是如何达到的。
	话虽如此，一旦你意识到该状态实际上是可达的，
	使用交互式工具找到该状态通常不会太难。
\item	这些工具对复杂的数据结构不太擅长，虽然
	可以使用形如 \qco{x=y; y=z; z=42;} 的
	初始化语句创建和遍历极其简单的链表。
\item	这些工具不处理内存映射 I/O 或设备寄存器。
	当然，处理这些东西需要将它们形式化，
	而这似乎不太可能。
\item	这些工具只会检测你为之编写 assertion 的问题。
	这个弱点是所有形式化方法所共有的，也是测试
	仍然重要的又一个原因。
	用本章开头引用的 Donald Knuth 的不朽名言来说，
	``当心上述代码中的 bug；
	我只是证明了它是正确的，而没有试过。''
\end{enumerate}

话虽如此，这些工具的一个优势是它们被设计为
模拟架构所允许的全部行为范围，包括
合法但当前硬件实现尚未施加于不知情的软件
开发者身上的行为。
因此，经过这些工具审查的算法在实际硬件上运行时
可能具有一些额外的安全边际。
此外，在真实硬件上的测试只能找到 bug；这样的测试
本质上无法证明给定的用法是正确的。
要理解这一点，考虑到研究人员在真实硬件上
常规运行超过1000亿次测试来验证他们的模型。
在一个案例中，尽管运行了1760亿次，
架构允许的行为没有出现~\cite{JadeAlglave2011ppcmem}。
相比之下，完整的 state-space search 允许工具
证明代码片段是正确的。

值得重复的是，形式化方法和工具不能替代测试。
事实是，生产大型可靠的并发软件工件（例如 Linux 内核）
是相当困难的。
因此，开发者必须准备好运用他们所掌握的每一个工具
来实现这个目标。
本章介绍的工具能够定位通过测试非常难以产生
（更不用说追踪）的 bug。
另一方面，测试可以应用于比本章介绍的工具
可能处理的大得多的软件体。
一如既往，使用合适的工具做合适的工作！

当然，最好是通过设计你的并行代码使其容易分区，
然后使用更高级别的原语（如锁、序列计数器、
原子操作和 RCU）来更直接地完成工作，
从而避免在这个层面工作的需要。
即使你绝对必须使用底层的内存屏障和读-修改-写指令
来完成工作，你越保守地使用这些锋利的工具，
你的生活就越可能轻松。
